\documentclass[11pt]{article}

\usepackage{amsmath, amssymb, amsthm}
\usepackage[margin=1in]{geometry}
\usepackage{enumitem}

\newtheorem{theorem}{Theorem}
\newtheorem{proposition}{Proposition}
\newtheorem{definition}{Definition}
\newtheorem{corollary}{Corollary}

\title{Mathematical Background: Topology Generator from a Subbasis}
\author{}
\date{}

\begin{document}
\maketitle

\section*{Scope}

$X$ is a fixed finite set throughout. This document covers: generating a topology from a
subbasis (\texttt{generator.py}); topology validation (\texttt{validate.py}); separation
axioms $T_0$--$T_4$ (\texttt{separation.py}); connectedness
(\texttt{connectedness.py}); compactness (\texttt{compactness.py}); and continuity
(\texttt{continuity.py}).

\section{Topologies, Bases, Subbases}

\begin{definition}
A \emph{topology} on $X$ is $\mathcal T \subseteq \mathcal P(X)$ with
$\emptyset, X \in \mathcal T$, closed under arbitrary unions and finite intersections.
\end{definition}

\begin{definition}
A \emph{basis} $\mathcal B$ generates $\mathcal T = \{\bigcup \mathcal U : \mathcal U
\subseteq \mathcal B\}$. A \emph{subbasis} $\mathcal S$ generates the basis
$\mathcal B = \{\bigcap \mathcal F : \mathcal F \subseteq \mathcal S,\ |\mathcal F| <
\infty\} \cup \{X\}$ (the empty intersection is $X$), and hence $\mathcal T$.
\end{definition}

Since $X$ is finite, $\mathcal P(X)$ is finite, so ``arbitrary union'' and ``finite union''
coincide, and every finite intersection is reachable by intersecting one subbasis element
at a time. \texttt{generate\_basis\_from\_subbasis} computes $\mathcal B$ by closing
$\{X\}$ under pairwise intersection with each $s \in \mathcal S$ to a fixed point;
\texttt{generate\_topology\_from\_basis} closes $\{\emptyset\}\cup\mathcal B$ under
pairwise union with each $b \in \mathcal B$ the same way. Both terminate since
$\mathcal P(X)$ is finite and each step either adds a new set or reaches the fixed point.

\begin{proposition}
For finite $X$, checking $\mathcal T$ closed under \emph{pairwise} union/intersection of
its own members is equivalent to closure under arbitrary unions and finite intersections.
\end{proposition}
\begin{proof}
Any union (intersection) of a subset of $\mathcal T$ is built by repeatedly combining two
elements at a time; if $\mathcal T$ is closed under each such step, by induction on
$|\mathcal T|<\infty$ it is closed under the full operation.
\end{proof}

This is what \texttt{is\_topology} checks directly, with a counterexample pair returned on
failure.

\section{Separation Axioms}

For distinct $x,y\in X$:

\begin{center}
\begin{tabular}{ll}
$T_0$ & some open set contains exactly one of $x,y$ \\
$T_1$ & some open set contains $x$ not $y$, and some open set contains $y$ not $x$ \\
$T_2$ & disjoint open sets separate $x$ and $y$ \\
\end{tabular}
\end{center}

\begin{proposition}
$T_1 \iff$ every singleton is closed.
\end{proposition}
\begin{proof}
($\Rightarrow$) For fixed $x$, $X\setminus\{x\} = \bigcup_{y\ne x} U_y$ where $U_y$ is open,
contains $y$, excludes $x$ (exists by $T_1$); a union of opens is open, so
$X\setminus\{x\}$ is open, i.e.\ $\{x\}$ is closed.
($\Leftarrow$) If every singleton is closed, $X\setminus\{y\}$ is an open set containing
$x$ but not $y$, and symmetrically for $X\setminus\{x\}$.
\end{proof}

\texttt{check\_T1\_via\_closed\_singletons} implements the right-hand side independently,
as a cross-check on \texttt{check\_T1}.

\begin{definition}[Munkres]
$X$ is \emph{regular} if for every closed $C$ and $x\notin C$, disjoint open sets separate
$\{x\}$ and $C$; \emph{normal} if every pair of disjoint closed sets is so separated.
$T_3 = \text{regular} \wedge T_1$; $T_4 = \text{normal}\wedge T_1$.
\end{definition}

\begin{theorem}
\label{thm:finite-t1-discrete}
Every finite $T_1$ space is discrete.
\end{theorem}
\begin{proof}
By the Proposition above, $T_1$ makes every singleton closed. $X$ finite $\Rightarrow$
every subset $A\subseteq X$ is a finite union of singletons, $A = \bigcup_{a\in A}\{a\}$;
each $X\setminus\{a\}$ is open (singletons closed), and $A = \bigcap_{a\notin A}
(X\setminus\{a\})$ is a finite intersection of opens, hence open. So every subset is open:
$\mathcal T = \mathcal P(X)$.
\end{proof}

\begin{corollary}
\label{cor:all-coincide}
On a finite space, $T_1 \iff T_2 \iff T_3 \iff T_4 \iff$ discrete.
\end{corollary}
\begin{proof}
$T_2\Rightarrow T_1$ always (fewer conditions). $T_1\Rightarrow$ discrete
(Theorem~\ref{thm:finite-t1-discrete}). Discrete $\Rightarrow T_2$ (singletons separate any
two points) $\Rightarrow$ discrete is regular and normal trivially (any two disjoint sets
are already separated by themselves, both open), so discrete $\wedge\, T_1 \Rightarrow
T_3, T_4$. $T_3, T_4 \Rightarrow T_1$ by definition. All five are therefore equivalent.
\end{proof}

Checked computationally across 500 random finite topologies in
\texttt{test\_t1\_t2\_t3\_t4\_and\_discrete\_all\_coincide\_on\_finite\_spaces}: for every
generated topology, \texttt{is\_T1 == is\_T2 == is\_T3 == is\_T4 == is\_discrete}.

\section{Connectedness}

\begin{definition}
$A\subseteq X$ is \emph{clopen} if $A$ and $X\setminus A$ are both open. $X$ is
\emph{connected} if its only clopen subsets are $\emptyset, X$.
\end{definition}

\begin{proposition}
The clopen subsets of a finite space form a Boolean algebra under $\cap,\cup,{}^c$, and for
$x\in X$, $\mathrm{comp}(x) := \bigcap\{C \text{ clopen} : x\in C\}$ is the connected
component of $x$.
\end{proposition}
\begin{proof}[Proof sketch]
Closure of clopen sets under $\cap,\cup,{}^c$ is immediate from the definition of a
topology. $\mathrm{comp}(x)$ is a finite intersection of clopen sets, hence clopen, and by
minimality no proper nonempty clopen subset of $\mathrm{comp}(x)$ can separate it further;
distinct $\mathrm{comp}(x)$'s are disjoint or equal since their intersection, if nonempty,
would itself be a smaller clopen set containing both would-be centers, contradicting
minimality unless they coincide.
\end{proof}

\texttt{connected\_components} computes $\mathrm{comp}(x)$ for each $x$ directly from this
characterization. The discrete topology: every singleton is clopen, so each
$\mathrm{comp}(x)=\{x\}$ ($n$ components). The indiscrete topology: only clopen sets are
$\emptyset,X$, so one component. The Sierpi\'nski space $(\{0,1\},
\{\emptyset,\{0\},\{0,1\}\})$: $\{0\}$ is open but $\{1\}$ (its complement) is not, so
$\{0\}$ is not clopen; the only clopen sets are $\emptyset, X$, so it is connected despite
not being $T_1$.

\section{Compactness}

\begin{definition}
$X$ is \emph{compact} if every open cover has a finite subcover.
\end{definition}

\begin{proposition}
Every finite topological space is compact.
\end{proposition}
\begin{proof}
$\mathcal T \subseteq \mathcal P(X)$ and $X$ finite $\Rightarrow |\mathcal T| \le
2^{|X|} < \infty$. Any open cover is a subset of $\mathcal T$, hence itself finite, hence
trivially its own finite subcover.
\end{proof}

\texttt{is\_compact} verifies this by constructing an explicit finite subcover via greedy
set cover (\texttt{has\_finite\_subcover}) rather than returning a hardcoded constant — the
content here is the proof that this is always possible, not the search itself.

\section{Continuity}

\begin{definition}
$f: (X,\mathcal T_X) \to (Y,\mathcal T_Y)$ is \emph{continuous} if $f^{-1}(V) \in
\mathcal T_X$ for every $V \in \mathcal T_Y$. A \emph{homeomorphism} is a continuous
bijection with continuous inverse.
\end{definition}

\begin{proposition}
If $\mathcal T_X$ is finer than $\mathcal T_X'$ (i.e.\ $\mathcal T_X' \subseteq
\mathcal T_X$), then $\mathrm{id}: (X,\mathcal T_X)\to(X,\mathcal T_X')$ is continuous.
\end{proposition}
\begin{proof}
$\mathrm{id}^{-1}(V) = V$ for every $V\in\mathcal T_X'\subseteq \mathcal T_X$.
\end{proof}

In particular $\mathrm{id}: \text{discrete} \to \text{any topology}$ is always continuous
(discrete is finest), and $\mathrm{id}:\text{any topology}\to\text{discrete}$ is continuous
only if the source is already discrete — confirmed directly:
\texttt{indiscrete}$\to$\texttt{discrete} fails, with a witnessing non-open preimage
returned by \texttt{is\_continuous}. A constant map is continuous regardless of
$\mathcal T_X$, since $f^{-1}(V) \in \{\emptyset, X\}$ for every $V$, and both are always
open. \texttt{is\_homeomorphism} checks $f,f^{-1}$ are genuine mutual-inverse bijections
before checking continuity in both directions.

\end{document}
